\subsection{Further complete linear-profile obstructions}

\begin{theorem}[Proved: a singleton long hom-set]
\label{CP2:A-baseline-singleton-long}
Let a three-object category have full linear acyclic support,
a singleton hom-set from the first to the last object, and
rationally acyclic corner classifying spaces.  Every positive
rational cohomology product of its classifying space is zero.
In particular this applies when all corner orders are at most
four.
\end{theorem}
\begin{proof}
Write $M,N,P$ for the corner monoids, $S,T$ for the two
adjacent bisets, and $U$ for the singleton long biset.
Decomposing nerve strings at their strict object changes
gives the triangular bar diagram with vertex spaces
$BM,BN,BP$, adjacent edge spaces
$A=B(N,S,M)$ and $B=B(P,T,N)$, long edge space
$K=B(P,U,M)=BP\times BM$, and face space
$F=B(P,T,N,S,M)$.

These are actual compatible bar maps.  A multi-simplex of
$F$ has the form $(p_*;t;n_*;s;m_*)$.  The maps to $A,B$
forget the opposite outer string and cross arrow.  The map
to $K$ retains $p_*,m_*$ and replaces the middle portion
by $t n_1\cdots n_q s$.  Middle faces preserve this
product; outer faces preserve it by the biset identities.
The composites to every vertex agree.  The homotopy
colimit of this diagram is $BC$: subdivision of the
products of these bar strings with their strict-object
simplex recovers the ordinary nerve.  In particular the
middle bar direction in $F$ is retained; no product
description $F=A\times B$ is assumed.

Replace the three vertex spaces and $K$ by points and
keep $A,B,F$ and both maps $F\to A,B$.  This is an actual
map of diagrams and an entrywise rational homology
equivalence.  For $K$ use rational Kunneth and the
endpoint acyclicity.  The finite homotopy-colimit
filtration makes the resulting map a rational homology
equivalence and hence a rational cohomology-ring
isomorphism.

Collapse the resulting long-edge interval, a contractible
CW subcomplex.  Its two endpoint vertices become one
suspension vertex; the middle vertex is the other.
Each triangular prism $F\times\Delta^2$ modulo the
long side is the suspension of an interval parametrized
by $F$, attached at its two ends through $F\to A,B$.
Thus the quotient is exactly the unreduced suspension
of the double mapping cylinder
$\operatorname{hocolim}(A\leftarrow F\to B)$.
All positive products vanish in that suspension, and
the collapse is a homotopy equivalence.  Pullback
through the rational comparison proves the theorem.
The final assertion uses the checked small-monoid
classification.
\end{proof}

\begin{corollary}[Proved: a projective initial singleton]
\label{CP2:A-baseline-projective-singleton}
Suppose a full-linear category has projective trivial left
rational module at its initial corner, a singleton first
adjacent hom-set, and final corner order at most nine.
Bounded rational
homology forces every rational degree-two product to
vanish.  In particular a profile $(a,b,c;x,y,z)$ with
$a\leq2$, $x=1$, and $c\leq9$ satisfies the hypothesis.
The opposite condition also suffices.
\end{corollary}
\begin{proof}
The initial corner's actual exact-orbit quotient replaces it
by a trivial object and replaces $x,z$ by their orbit
sets, inducing a rational homology and cohomology-ring
isomorphism.  The first adjacent hom-set remains a
singleton.  The comma category for deleting the middle
object is therefore a point: there is just the unique
incoming arrow from the now trivial first corner.
The deletion is an actual homotopy equivalence to
the ordinary left-act cone on the remaining long
hom-set over the final corner.  Its boundedness
follows from these comparisons.  The bounded-nine
act theorem, Theorem~\ref{CP2:A-bounded-nine-acts},
annihilates its second products.  Pull back and take
opposites for the other orientation.
\end{proof}

\begin{proposition}[Proved: two small incoming stalks]
\label{CP2:A-baseline-two-small-incoming}
Let a full-linear three-object category have trivial
first two corners, a last corner $M$ of order at most
eight, two arrows from the middle to the last object,
and at most three arrows from the first to the last.
It cannot realize $\mathbb{CP}^2$.  There is no bound
on the first adjacent hom-set.
\end{proposition}
\begin{proof}
Write $X=C(0,1)$, $Y=C(1,2)$ and $Z=C(0,2)$.
The first two objects form a graph category $D$, so
$H^2(BD;\mathbb Q)=0$.  If $H^2(BM;\mathbb Q)=0$,
the two-block argument applies.  Otherwise the small
cardinality and the minimal-ideal row and column
requirements force trivial minimal group.  The
universal group of $M$ is then trivial: for a minimal
idempotent $e$, every group image satisfies
$\eta(m)=\eta(eme)$, and $eMe=\{e\}$.

Each $x\in X$ gives an $M$-equivariant map
$\beta_x:Y\to Z$.  A constant such map supplies a
globally $M$-fixed point of $Z$, so the fixed-arrow
corollary excludes all second products.  Assume all
these maps are injective.  If $|Z|=3$ and two image
sets differ, their singleton intersection is invariant
under $M$ and again gives a fixed point.  Hence all
images are one common two-element invariant set $U$.
The same holds automatically when $|Z|=2$.

Choose one bijection $Y\to U$ to identify these two
sets.  If two of the maps $\beta_x$ differ, a
transposition commutes with every actor on $Y$.
Its transformation centralizer on two points consists
of permutations.  The resulting homomorphism from
$M$ to that permutation group is trivial because
$U(M)=1$.  Thus every point of $U$ is fixed, and
the fixed-arrow obstruction applies.  The same
conclusion holds if all maps agree and the $M$ action
on $Y$ consists of permutations.

In the only remaining case all $\beta_x$ agree and
some actor on $Y$ is constant.  The action category
of $Y$ is connected, since that constant joins both
states to its value.  Collapse $X$ to a singleton,
obtaining an actual category $C_0$ and functor
$q:C\to C_0$.  Choosing one member of $X$ gives a
section $j$; all compositions are preserved because
the maps $\beta_x$ coincide.

Here is the exact second-cohomology comparison.
Relative to the section category, a normalized chain
uses one extra arrow of $X$, optionally followed by
an arrow of $Y$ and a string in $M$.  For each extra
arrow the relative chain complex is the augmentation
cone of the chains of the $Y$ action category, shifted
once.  Its degree-two homology is
$\widetilde H_0(Y_{hM};\mathbb Q)=0$.  Consequently
$H^2(BC;\mathbb Q)\to H^2(BC_0;\mathbb Q)$ is
injective, and the retraction makes it an isomorphism
with inverse $q^*$.  This argument does not assume
separate boundedness of any corner or action space.

Under a putative projective-plane hypothesis, $BC_0$
has bounded rational homology because it is an actual
retract.  Its first adjacent hom-set is singleton
with trivial source, so deleting the middle object
is an actual homotopy equivalence to an ordinary
act cone over $M$ with $|Z|\leq3$ states.  The
bounded order-eight act theorem kills all its
degree-two products.  The displayed isomorphism
then kills the square on $BC$, a contradiction.
The case $|Z|=1$ was already the fixed-arrow case.
\end{proof}
